% Учебный исследовательский проект / Research project.
% Build: xelatex main; bibtex main; xelatex main; xelatex main.
% Style from https://github.com/kisnikser/m1p-template/tree/main/paper
\documentclass{article}
\PassOptionsToPackage{numbers,sort,compress}{natbib}
\usepackage[main,final]{neurips_2025}
\usepackage{fontspec}
\usepackage[english,russian]{babel}
\setmainfont{cmunrm.otf}[BoldFont=cmunbx.otf,ItalicFont=cmunti.otf,BoldItalicFont=cmunbi.otf]
\setsansfont{cmunss.otf}
\setmonofont{cmuntt.otf}
\usepackage[hidelinks]{hyperref}
\usepackage{url}
\usepackage{booktabs}
\usepackage{amsmath,amssymb}
\makeatletter
\renewcommand{\@noticestring}{Учебный исследовательский проект / Research project.}
\makeatother
\renewenvironment{abstract}{\vskip 0.075in\centerline{\large\bfseries\abstractname}\vspace{0.5ex}\begin{quote}}{\par\end{quote}\vskip 1ex}
\setlength{\emergencystretch}{2em}
\title{Адаптивное условное диффузионное сэмплирование при поступлении новых наблюдений в процессе генерации\\[0.5em]{\large Adaptive Conditional Diffusion Sampling with Observations Arriving During Generation}}
\author{}
\begin{document}
\maketitle
\begin{abstract}
Данная работа посвящена исследованию способов адаптации диффузионных моделей для генерации траекторий к условиям, когда новый контекст поступает во время генерации. В задачах прогнозирования движения дополнительные наблюдения могут уточнить возможные варианты будущего или существенно изменить их вероятности, однако учёт этих данных может потребовать повторных вычислений, которые не всегда допустимы при ограниченном времени и вычислительном бюджете. Я исследую, как обновлять формируемую траекторию с учётом новых данных, в каких случаях можно использовать уже выполненные вычисления и как выбранный способ обновления влияет на качество прогноза и затраты на его получение. Для этого я рассмотрю существующие подходы к диффузионному прогнозированию, потоковой генерации и условному сэмплированию, сопоставлю их предположения и ограничения и выделю способы, применимые к изменению контекста внутри незавершённого процесса генерации. Сравнение предполагается проводить при разных моментах поступления наблюдений, разной степени их влияния на прогноз и сопоставимых вычислительных бюджетах, учитывая точность траекторий, разнообразие возможных исходов и время получения обновлённого результата. При этом я отдельно рассмотрю качество вероятностного прогноза, поскольку близость одной сгенерированной траектории к наблюдаемому движению ещё не означает, что модель корректно отражает неопределённость будущего. На основе выявленных ограничений я предложу возможные способы улучшения существующих подходов и проверю, в каких условиях они позволяют эффективнее использовать новые наблюдения. Исследование направлено на определение возможностей и границ адаптации диффузионного сэмплирования к изменяющемуся контексту при ограниченных вычислительных ресурсах.
\end{abstract}
\noindent\textbf{Ключевые слова:} диффузионные модели; генерация траекторий; вероятностное прогнозирование; условное сэмплирование; потоковые наблюдения; вычислительная эффективность.

\clearpage
\begin{otherlanguage}{english}
\begin{abstract}
This work investigates ways to adapt diffusion models for trajectory generation to settings where new context arrives during generation. In motion forecasting, additional observations can refine possible futures or substantially change their probabilities, but incorporating these data may require repeated computation that is not always feasible under limited time and computational budgets. I investigate how to update a trajectory being generated using new data, when previous computation can be reused, and how the update strategy affects forecast quality and computational cost. To this end, I will review existing approaches to diffusion forecasting, streaming generation, and conditional sampling, compare their assumptions and limitations, and identify strategies applicable to context changes within an unfinished generation process. The comparison will consider different observation arrival times, different degrees of influence on the forecast, and comparable computational budgets, assessing trajectory accuracy, diversity of possible outcomes, and the time required to produce an updated result. I will also examine probabilistic forecast quality, since a single generated trajectory being close to the observed motion does not necessarily mean that the model correctly represents uncertainty about the future. Based on the identified limitations, I will propose possible improvements and test the conditions under which they make more effective use of new observations. The investigation aims to establish the capabilities and limits of adapting diffusion sampling to changing context under limited computational resources.
\end{abstract}
\noindent\textbf{Keywords:} diffusion models; trajectory generation; probabilistic forecasting; conditional sampling; streaming observations; computational efficiency.
\end{otherlanguage}

\clearpage
\section{Обзор литературы}
\subsection{Диффузионная генерация и прогнозирование траекторий}

Диффузионные модели формируют результат через последовательное расшумление. Базовая дискретная формулировка представлена в DDPM \citep{p37}, а непрерывная связь между стохастическим процессом и дифференциальным уравнением — в Score-SDE \citep{p38}. Для моей темы важно, что между начальным шумом и готовой траекторией существует промежуточное состояние. При поступлении новых данных возникает вопрос, можно ли использовать это состояние для обновлённого прогноза и какие ошибки при этом появляются.

В прогнозировании движения MID использует историю траекторий и взаимодействия участников для условной генерации \citep{p09}. LED дополнительно обучает инициализацию, позволяющую пропускать часть расшумления \citep{p10}, а MotionDiffuser моделирует совместные варианты движения нескольких участников и поддерживает направленное сэмплирование \citep{p11}. Эти работы дают предметную основу для сравнения качества, разнообразия и стоимости прогнозов. DiffTraj также работает с траекториями, но решает задачу генерации синтетических GPS-данных \citep{p12}. Поэтому я рассматриваю её как источник идей о представлении данных, а не как прямой аналог обновления прогноза.

Отдельную группу составляют методы планирования. Diffuser и Decision Diffuser используют условную генерацию для построения поведения с заданными целями и ограничениями \citep{p14}, \citep{p15}. Diffusion Policy генерирует последовательности действий по наблюдениям \citep{p17}; Diffusion Planner рассматривает автономное движение \citep{p13}, а DiffPhyCon — управление физическими системами \citep{p18}. Здесь качество определяется не только соответствием распределению данных, но и результатом исполнения выбранных действий. AdaptDiffuser дополнительно изменяет сам планировщик через дообучение на отобранных синтетических данных \citep{p16}. Я отделяю такую адаптацию модели от обновления одного незавершённого прогноза.

\subsection{Новые наблюдения и повторное использование вычислений}

Наиболее близки к теме работы, сохраняющие часть генерации между поступлениями данных. Streaming Diffusion Policy поддерживает буфер действий с разными уровнями шума и уточняет его по новым наблюдениям \citep{p01}. CL-DiffPhyCon также использует асинхронное расшумление, включая обратную связь при генерации дальнейших управляющих воздействий \citep{p02}. Они показывают, что использование нового контекста до завершения всей последовательности уже является самостоятельным направлением исследований. Diffusion Forcing даёт более общий механизм обучения последовательностей с разными уровнями шума у элементов \citep{p05}.

RTI-DP использует ранее предсказанную последовательность как начальное приближение для следующего расчёта \citep{p04}. Sequential Flow Matching рассматривает перенос предыдущего распределения в обновлённое по наблюдениям \citep{p03}. В отличие от повторного использования отдельной траектории, здесь обучается переход между распределениями. Авторы последней работы отмечают накопление ошибок при рекурсивных обновлениях и обсуждают адаптивное повторное зашумление и перезапуск как дальнейшие направления. Это конкретная связь с моей темой, но не основание считать такую общую идею новой.

Задержку можно учитывать и на уровне исполнения. RTC согласует соседние фрагменты действий при асинхронной генерации \citep{p06}. Reactive Diffusion Policy выделяет быстрый контур обратной связи \citep{p07}, а FA-RDP адаптирует частоту вывода и число шагов генерации к неоднозначности действий \citep{p08}. Эти подходы показывают, что реакция на новые данные может достигаться разными механизмами. Поэтому я буду различать обновление промежуточного сэмпла, перенос готового прогноза и коррекцию исполнения: близкие практические цели не делают эти задачи одинаковыми.

\subsection{Условное распределение по наблюдениям}

Для анализа обновлений полезны работы с временными рядами. CSDI и SSSD восстанавливают пропущенные значения по доступным наблюдениям \citep{p19}, \citep{p23}. TimeGrad строит авторегрессионный прогноз \citep{p20}, а TimeDiff генерирует будущий участок совместно и использует специальные механизмы задания условия \citep{p21}. TSDiff вводит условия при выводе и рассматривает уточнение готовых прогнозов \citep{p24}. Эти методы помогают выбрать основу для экспериментов, но применение к наблюдениям, поступающим внутри расшумления, требует собственного протокола.

Score-based Data Assimilation отделяет модель наблюдений от обучения генератора и восстанавливает траектории по неполным или шумным измерениям \citep{p22}. В обратных задачах похожую роль играют DDRM и Diffusion Posterior Sampling \citep{p29}, \citep{p28}. Однако согласование с измерениями может опираться на приближения, а ограничения линейной модели наблюдений подходят не для любого контекста. Кроме того, Classifier-Free Guidance управляет соответствием условию и разнообразием \citep{p40}; усиление такого соответствия само по себе не подтверждает корректность вероятностного обновления.

\subsection{Вычислительная эффективность и повторное зашумление}

Сокращение затрат не обязательно связано с адаптацией контекста. DDIM, DPM-Solver, DPM-Solver++ и UniPC ускоряют численное сэмплирование \citep{p30}, \citep{p31}, \citep{p32}, \citep{p33}. EDM систематизирует выбор параметризации, расписания шума и сэмплера \citep{p36}. Progressive Distillation и Consistency Models уменьшают число шагов через изменение обучения \citep{p35}, \citep{p34}, а Flow Matching представляет смежный способ обучать перенос распределений \citep{p39}. Эти работы нужны, чтобы отделить возможную пользу обновления контекста от обычного ускорения генератора и учитывать стоимость дополнительного обучения.

Повторное зашумление также является известным инструментом. RePaint использует повторное сэмплирование при восстановлении изображений \citep{p27}, Restart исследует сочетание добавления шума и обратного интегрирования \citep{p25}, а RePS переносит restart-подход на обратные задачи \citep{p26}. Их результаты обосновывают включение подобных операций в сравнение, но не устанавливают, какая коррекция нужна после конкретного нового наблюдения в задаче траекторий.

\subsection{Направление исследования}

Рассмотренные работы уже предлагают способы потоковой генерации, повторного использования прогнозов и ускорения сэмплирования. Поэтому я не формулирую задачу как создание первого метода, учитывающего новые наблюдения. Я исследую, как эти способы работают при изменении контекста внутри незавершённой генерации траектории и как их применимость зависит от момента поступления данных, величины изменения прогноза и доступного бюджета. Для сравнения важны и точность, и сохранение неопределённости, и время получения обновления. На основе выявленных ограничений я предложу возможные улучшения; конкретный подход и его отличие от ближайших методов будут определены по результатам дальнейшего анализа и экспериментов.

\clearpage
\begin{otherlanguage}{english}
\section{Related Work}
\subsection{Diffusion generation and trajectory forecasting}

Diffusion models generate outputs through successive denoising. DDPM provides the basic discrete formulation \citep{p37}, while Score-SDE relates stochastic processes and differential equations in continuous time \citep{p38}. For my topic, an important feature is the intermediate state between initial noise and a completed trajectory. When new data arrive, the question is whether this state can support an updated forecast and what errors this introduces.

In motion forecasting, MID conditions generation on trajectory history and interactions \citep{p09}. LED additionally learns an initialization that skips part of denoising \citep{p10}, while MotionDiffuser models joint multi-agent futures and supports guided sampling \citep{p11}. These studies provide a domain-specific basis for comparing forecast quality, diversity, and cost. DiffTraj also addresses trajectories but generates synthetic GPS data \citep{p12}. I therefore consider it a source of representation ideas rather than a direct analogue of forecast updating.

Planning methods form a separate group. Diffuser and Decision Diffuser use conditional generation to construct behavior with specified goals and constraints \citep{p14}, \citep{p15}. Diffusion Policy generates action sequences from observations \citep{p17}; Diffusion Planner addresses autonomous driving \citep{p13}, and DiffPhyCon addresses physical-system control \citep{p18}. Their quality depends on the outcome of executed actions as well as agreement with the data distribution. AdaptDiffuser additionally changes the planner through fine-tuning on selected synthetic data \citep{p16}. I distinguish this model adaptation from updating a single unfinished forecast.

\subsection{New observations and reuse of computation}

The closest studies preserve part of generation across incoming observations. Streaming Diffusion Policy maintains an action buffer with different noise levels and refines it using new observations \citep{p01}. CL-DiffPhyCon similarly uses asynchronous denoising, incorporating feedback when generating subsequent controls \citep{p02}. They demonstrate that using new context before the full sequence is complete is already an established research direction. Diffusion Forcing provides a more general training mechanism for sequences with different noise levels per element \citep{p05}.

RTI-DP uses a previously predicted sequence as the initial guess for the next computation \citep{p04}. Sequential Flow Matching transports the previous distribution to an updated one conditioned on observations \citep{p03}. Unlike reusing an individual trajectory, it learns a transition between distributions. Its authors identify error accumulation in recursive updates and discuss adaptive re-noising and resets as future directions. This directly connects to my topic but does not make the general idea novel.

Latency can also be addressed at execution time. RTC aligns adjacent action chunks under asynchronous generation \citep{p06}. Reactive Diffusion Policy introduces a fast feedback loop \citep{p07}, while FA-RDP adapts inference frequency and sampling steps to action ambiguity \citep{p08}. These approaches demonstrate different mechanisms for responding to new data. I will therefore distinguish intermediate-sample updates, reuse of completed forecasts, and execution corrections: similar practical goals do not make these problems identical.

\clearpage
\subsection{Conditional distributions from observations}

Time-series studies provide useful foundations for analyzing updates. CSDI and SSSD impute missing values from available observations \citep{p19}, \citep{p23}. TimeGrad produces autoregressive forecasts \citep{p20}, while TimeDiff generates the future segment jointly with specialized conditioning mechanisms \citep{p21}. TSDiff introduces conditions at inference and considers refining existing forecasts \citep{p24}. These methods inform the choice of an experimental base, but observations arriving within denoising require a dedicated protocol.

Score-based Data Assimilation separates the observation model from generator training and reconstructs trajectories from incomplete or noisy measurements \citep{p22}. DDRM and Diffusion Posterior Sampling play related roles in inverse problems \citep{p29}, \citep{p28}. However, measurement conditioning can rely on approximations, and linear observation assumptions do not cover every form of context. Classifier-Free Guidance controls condition adherence and diversity \citep{p40}; increasing adherence does not itself establish a correct probabilistic update.

\subsection{Computational efficiency and re-noising}

Reducing cost does not necessarily require context adaptation. DDIM, DPM-Solver, DPM-Solver++, and UniPC accelerate numerical sampling \citep{p30}, \citep{p31}, \citep{p32}, \citep{p33}. EDM organizes choices of parameterization, noise scheduling, and sampling \citep{p36}. Progressive Distillation and Consistency Models reduce sampling steps through changes in training \citep{p35}, \citep{p34}, while Flow Matching provides an adjacent approach to learning distribution transport \citep{p39}. These studies help separate the potential benefit of context updating from ordinary generator acceleration and account for additional training costs.

Re-noising is also an existing tool. RePaint uses resampling for image reconstruction \citep{p27}, Restart studies adding noise followed by reverse integration \citep{p25}, and RePS extends restart-based sampling to inverse problems \citep{p26}. Their results motivate including such operations in comparisons but do not establish which correction is appropriate after a particular new observation in trajectory forecasting.

\subsection{Research direction}

The reviewed studies already offer streaming generation, forecast reuse, and accelerated sampling. I therefore do not frame the problem as creating the first method that incorporates new observations. I investigate how these strategies behave when context changes within unfinished trajectory generation and how their applicability depends on observation arrival time, the extent of forecast change, and the available budget. The comparison must consider accuracy, uncertainty preservation, and update latency. Based on identified limitations, I will propose possible improvements; the particular approach and its distinction from the closest methods will be determined through further analysis and experiments.
\end{otherlanguage}

\clearpage
\bibliographystyle{unsrtnat}
\bibliography{references}
\end{document}
